\documentclass[11point]{article}
% set the margins
\usepackage[left=25mm,right=25mm,top=20mm,bottom=20mm]{geometry}
% for displaying images
\usepackage{graphicx}
% for equation support, e.g. eqref
\usepackage{amsmath}
% for table support, e.g. toprule
\usepackage{booktabs}
% for multiple authors as a block
\usepackage{authblk}

% I redefine vectors to be bold
\renewcommand{\vec}[1]{\boldsymbol{\mathbf{#1}}}

% I can define new macros for things I use a lot
\newcommand{\datapoint}{\vec{x}}
\newcommand{\weights}{\vec{w}}

\DeclareMathOperator{\argmax}{argmax}
\twocolumn

\begin{document}

\title{Report Title}
\author[*]{Author Name1}
\author[**]{Author Name2}
\affil[*]{Dept. of Information Studies}
\affil[**]{Dept. of Physics \& Astronomy}
\affil[ ]{University College London, WC1E 6BT}
\affil[ ]{\textit {\{email1,email2\}@ucl.ac.uk}}
\date{\today}

\maketitle

\begin{abstract}
The abstract text goes here. Here is a bit more of my abstract.
\end{abstract}

\section{Introduction}
\label{sec:intro}
Here is the text of your introduction.

Here is a simple equation
% note this uses the equation environment
\begin{equation}
    \label{eqn:simple_equation}
    \alpha = \sqrt{ \beta }
\end{equation}

I can refer to my equation by saying look at Equation \eqref{eqn:simple_equation}.

Here is a multi-line equation:
% note this uses the align environment
\begin{align}
% \input and \weights in this equation are defined using a macro 
% see the \newcommand lines at the top
f(\datapoint; \weights)
% & helps neatly align equations over multiple lines 
& = \sigma(\weights^T\datapoint)
\label{eqn:logistic_regresion_decision_function} % label
\\ % line break
& = \frac{1}{1+e^{-\weights^T\datapoint}}
\notag % don't give this line an equation number
\end{align}

I can point people to literature like this \cite{Mikolov2013} (a conference paper) or this \cite{Fawcett2006} (a journal paper), or even this \cite[Sec. 4.1]{Bishop2006} (a book, with section indicated). See the file \texttt{report.bib} for how these bibliographic entries are defined.

It helps to have a convention to the font used for particular mathematical objects. A common convention is: for simple values use lower case latin characters $s$ (this might include observed values in probabilistic work); for sets use caligraphic upper case $\mathcal{S}$; for vectors of values use bold latin $\vec{s}$; for matrices use upper case latin $S$; and for parameters (including parameters in probabilistic models) use lower case greek letters (scalar), $\sigma$, bold greek (vectors), $\vec{\sigma}$, and possibly upper-case greek for matrices of parameters $\Sigma$. However, it is up to you to develop your own style.

\subsection{Subsection Heading}
Write your subsection text here. Here I am referencing Section \ref{sec:intro}. Figures and tables can be inserted and referred to, e.g. Figure \ref{fig:simulation_figure} and Table \ref{tbl:results_table}.

\begin{figure}
    \centering
    \includegraphics[width=3.0in]{myfigure}
    \caption{Simulation Results}
    \label{fig:simulation_figure}
\end{figure}

\begin{table}[htb]
  \centering
  \begin{tabular}{lcc}
  \toprule
  & Interested & Understand \\
  \midrule
  Linear Regression & N & Y \\
  Neural Networks & Y & N \\
  \bottomrule
  \end{tabular}
  \caption{Some results as a table.}
  \label{tbl:results_table}
\end{table}

\subsection{Quoting and Emphasis}
When you quote or emphasise things, you have a few options. I tend to put things in italics, e.g. Luke told me, \emph{Put it in italics!}. If you prefer to put it in quote-marks then do so, but make sure you use left and right quotes properly, e.g. Luke told me, ``Use left quotes to open, and right quotes to close!''.


% sections for your report can be in subfiles
\section{Background}
Your background section text goes here\ldots

\section{Method}
Your method section text goes here\ldots

\section{Results}
Your results section text goes here\ldots

\section{Discussion}
Your discussion section text goes here\ldots



\section*{Declaration}
This document represents the group report submission from the named authors for the project assignment of module: Foundations of Machine Learning and Data Science (INST0060), 2019-20. In submitting this document, the authors certify that all submissions for this project are a fair representation of their own work and satisfy the UCL regulations on  plagiarism.

% to put references in
\bibliography{report}
% define the bibliography style
\bibliographystyle{plain}

\end{document}
